\documentclass[11pt]{article}

\usepackage[margin=1in]{geometry}
\usepackage{amsmath, amssymb, amsthm}
\usepackage{mathtools}
\usepackage{bm}
\usepackage{xcolor}
\usepackage{hyperref}
\usepackage[most]{tcolorbox}
\usepackage{booktabs}
\usepackage{enumitem}
\usepackage{listings}

\hypersetup{colorlinks=true, linkcolor=blue!50!black, urlcolor=blue!50!black,
            unicode=true, bookmarksdepth=2}

\definecolor{codebg}{HTML}{F3F4F6}
\definecolor{accent}{HTML}{C85020}
\definecolor{navy}{HTML}{12264E}

\lstset{
  basicstyle=\ttfamily\small,
  backgroundcolor=\color{codebg},
  frame=none, breaklines=true, columns=fullflexible,
  keepspaces=true, xleftmargin=4pt, framexleftmargin=4pt,
}

\newcommand{\code}[1]{\texttt{\small #1}}
\newcommand{\Adoref}[1]{\textcolor{accent}{\footnotesize\textbf{[ado L#1]}}}

\newcommand{\R}{\mathbb{R}}
\newcommand{\E}{\mathbb{E}}
\newcommand{\1}{\mathbb{1}}
\newcommand{\Smpl}{\mathcal{S}}
\newcommand{\thetad}{\widehat{\bm{\theta}}_d}
\newcommand{\Ngt}{N_{g,t}}

\title{\textbf{Stages 2 and 3 of \texttt{did\_multiplegt\_dyn}}\\[2pt]
       \large Where Stage 1's controls outputs are consumed: long-difference
       residualization and the $U^{\mathrm{var},X}$ variance correction}
\author{Notes on L4582--L4954 of \texttt{did\_multiplegt\_dyn.ado}}
\date{}

\begin{document}
\maketitle

\noindent
\begin{tcolorbox}[colback=codebg, colframe=navy, boxrule=0.6pt, arc=1pt]
Stage 1 (the previous note) produced, for each baseline level $d$,
\begin{itemize}[leftmargin=*, itemsep=0pt]
\item $\thetad \in \R^{K}$ \,---\, \code{coefs\_sq\_`d'\_XX};
\item $\mathrm{Den}_d^{-1} \in \R^{K\times K}$ \,---\, \code{inv\_Denom\_`d'\_XX};
\item the per-control score series $\mathrm{prod\_X}_{k,g,t}$ \,---\, \code{prod\_X{k}\_Ngt\_XX};
\item the cell-mean fitted value $\widehat{E}^{(d)}_{g,t}$ \,---\, \code{E\_y\_hat\_gt\_int\_`d'\_XX}.
\end{itemize}
This note walks through how Stages 2 and 3 spend each of those four objects.
\textbf{Stage 2} replaces the long-difference of $Y$ with a controls-residualized
version; \textbf{Stage 3} adds a variance term that accounts for the fact that
$\thetad$ is estimated, not known. Effects-side only; placebos are symmetric and
covered briefly at the end.
\end{tcolorbox}

\section{Notation refresh}

\noindent
\begin{tabular}{@{}p{4.6cm}p{10.4cm}@{}}
\toprule
$\ell \in \{1,\dots,L\}$           & event-study horizon (\code{`i'} in the loop). \\
$d \in \{1,\dots,L_D\}$            & baseline treatment level (\code{d\_sq\_int\_XX}).\\
$F_g$                              & first switch date for group $g$ (\code{F\_g\_XX}); $F_g=T+1$ for never-switchers.\\
$T_g$                              & last period where $g$ still has a control group (\code{T\_g\_XX}).\\
$S_g \in \{0,1\}$                  & switcher direction: 1 = switcher-in ($\pm = +$), 0 = switcher-out ($\pm = -$).\\
$D_{g,1}$                          & baseline treatment of group $g$.\\
$Y_{g,t}, X_{k,g,t}$               & outcome and $k$-th control.\\
$\Delta Y^{(\ell)}_{g,t} = Y_{g,t} - Y_{g,t-\ell}$ & long-difference of $Y$, \code{diff\_y\_`l'\_XX}.\\
$\Delta X^{(\ell)}_{k,g,t} = X_{k,g,t} - X_{k,g,t-\ell}$ & long-difference of $X_k$, \code{diff\_X{k}\_`l'\_XX}.\\
$\Ngt$                             & weight of cell $(g,t)$ (\code{N\_gt\_XX}).\\
$G$                                & number of groups (\code{G\_XX}).\\
$N^{\pm}_{\ell}$                   & number of switcher-$\pm$ groups eligible at horizon $\ell$ (\code{N`+\!-'\_`l'\_XX}).\\
\bottomrule
\end{tabular}

\medskip\noindent
Throughout, $\pm$ stands for ``$+$ for switchers-in (\code{increase\_XX==1}),
$-$ for switchers-out (\code{increase\_XX==0})''. The code generates both
families by setting \code{increase\_XX} in an outer loop, so the same lines
build $+$ and $-$ versions in turn.

% ============================================================================
\section[Stage 2 -- Long-difference residualization (L4582--4689)]{Stage 2 \,---\, Long-difference residualization \Adoref{4582--4689}}

\subsection*{Step 2.1 \,---\, Build the long differences of each control}

For each horizon $\ell$ (looping over $\ell=1,\dots,L$ outside this block) and
each control $k$ \Adoref{4586--4598}:
\begin{align}
\Delta X^{(\ell)}_{k,g,t} \;&=\; X_{k,g,t} - X_{k,g,t-\ell}, \label{eq:longDX}\\
\Delta X^{(\ell),N}_{k,g,t} \;&=\; \Ngt\cdot \Delta X^{(\ell)}_{k,g,t}. \label{eq:longDXN}
\end{align}
Stored as \code{diff\_X{k}\_`l'\_XX} and \code{diff\_X{k}\_`l'\_N\_XX}.

\subsection*{Step 2.2 \,---\, The treatment-side moment $M^{\pm}_{d,\ell,k}$ \Adoref{4611--4622}}

The companion-paper vectors $\bm{m}^{\pm}_{g,d,\ell}$ have one coordinate per
control $k$. The code builds them in three passes.

\medskip
\textbf{Per-$(g,t)$ summand} \Adoref{4612}:
\begin{equation}
m^{\pm}_{g,d,\ell,k,t} \;=\; \1\{\ell\le T_g - 2,\ D_{g,1}=d\}\cdot \frac{G}{N^{\pm}_{\ell}}\cdot
\Bigl[\,\mathrm{dist}^{(\ell)}_{g,t} - \tfrac{N^{\pm,t}_{\ell,g}}{N^{\mathrm{c},(\ell)}_{g,t}}\cdot\mathrm{nc}^{(\ell)}_{g,t}\,\Bigr]\cdot \1\{\ell{+}1\le t\le T_g\}\cdot \Delta X^{(\ell),N}_{k,g,t}.
\label{eq:msummand}
\end{equation}
The four named quantities (all built earlier in the file) are:
\begin{itemize}[leftmargin=*, itemsep=0pt]
\item $\mathrm{dist}^{(\ell)}_{g,t}$ \,---\, distance-to-switch indicator
       (\code{distance\_to\_switch\_`l'\_XX}): equals 1 when $(g,t)$ is at
       relative time $\ell$ for switcher $g$.
\item $\mathrm{nc}^{(\ell)}_{g,t}$ \,---\, never-changer indicator
       (\code{never\_change\_d\_`l'\_XX}): 1 if $g$ is still untreated at $t$
       (control side).
\item $N^{\pm,t}_{\ell,g}$ \,---\, cohort-size count
       (\code{N`+\!-'\_t\_`l'\_g\_XX}).
\item $N^{\mathrm{c},(\ell)}_{g,t}$ \,---\, count of controls available for $g$
       at horizon $\ell$ (\code{N\_gt\_control\_`l'\_XX}).
\end{itemize}

\textbf{Sum across $t$ within $g$} \Adoref{4615}:
\begin{equation}
m^{\pm}_{g,d,\ell,k} \;=\; \sum_{t} m^{\pm}_{g,d,\ell,k,t},
\label{eq:msumt}
\end{equation}
kept on the first row of each group (\code{... = . if \_n != 1}, L4616).

\textbf{Average over groups} \Adoref{4620--4622}:
\begin{equation}
\boxed{\;M^{\pm}_{d,\ell,k} \;=\; \frac{1}{G}\sum_{g} m^{\pm}_{g,d,\ell,k}.\;}
\label{eq:Mdlk}
\end{equation}
Stored as \code{M`+\!-'\_`d'\_`k'\_`l'\_XX}.

\medskip\noindent
\textbf{What this is.} $M^{\pm}_{d,\ell,k}$ is the cross-group average of the
$k$-th control's long-difference, weighted by the residual of
the distance-to-switch indicator on its control-group mean. It is the
``treatment-side'' summary that the variance correction multiplies
$(\mathrm{Den}_d^{-1}\cdot\mathrm{score}_d - \thetad)$ by.

\subsection*{Step 2.3 \,---\, The cell-mean fitted value $\widehat{E}^{(d)}_{g,t}$ \Adoref{4625--4647}}

The Stage-1 fitted value $\widehat{E}^{(d)}_{g,t,\mathrm{int}} =$
\code{E\_y\_hat\_gt\_int\_`d'\_XX} is now masked to zero in
cells with too few control groups:
\begin{equation}
\widehat{E}^{(d,\ell)}_{g,t} \;=\; \widehat{E}^{(d)}_{g,t,\mathrm{int}}\cdot \1\!\bigl\{E^{\mathrm{denom}}_{d,\ell}(t) \ge 2\bigr\},
\label{eq:Ehat}
\end{equation}
where $E^{\mathrm{denom}}_{d,\ell}(t)$ counts the available not-yet-switcher
control groups (or clusters, if \code{cluster()} is used) at time $t$ inside
baseline $d$ \Adoref{4625--4641}. Stored as \code{E\_y\_hat\_gt\_`d'\_XX}.

\subsection*{Step 2.4 \,---\, The per-group control-side score $\mathrm{in\_sum}_{k,d}$ \Adoref{4653--4669}}

Two scalar building blocks. First, the within-baseline window-size weight
\Adoref{4657--4659}:
\begin{equation}
N^{\mathrm{c}}_d \;=\; \sum_{(g,t)} \Ngt\cdot \1\!\left\{D_{g,1}=d,\ 2\le t\le T_d,\ t<F_g,\ \Delta Y_{g,t}\ne\text{miss}\right\},
\label{eq:Nc}
\end{equation}
where $T_d = \max_{h:D_{h,1}=d} F_h - 1$ \Adoref{4578--4580}.

\medskip
Second, the per-$(g,t)$ score summand \Adoref{4664}:
\begin{equation}
\sigma^{(d)}_{k,g,t} \;=\; \frac{\mathrm{prod\_X}_{k,g,t}\cdot \delta^{(d)}_{g,t}\cdot \bigl(\Delta Y_{g,t} - \widehat{E}^{(d,\ell)}_{g,t}\bigr)\cdot \1\{2\le t\le F_g{-}1\}}{N^{\mathrm{c}}_d},
\label{eq:sigma}
\end{equation}
with the DoF adjustment
\begin{equation}
\delta^{(d)}_{g,t} \;=\; 1 + \1\!\bigl\{E^{\mathrm{denom}}_{d,\ell}(t) \ge 2\bigr\}\cdot
\left[\sqrt{\tfrac{E^{\mathrm{denom}}_{d,\ell}(t)}{E^{\mathrm{denom}}_{d,\ell}(t)-1}}-1\right]
\label{eq:dof}
\end{equation}
\Adoref{4661--4663}. The adjustment is the standard small-sample correction
when demeaning by an estimated cell mean.

\medskip
Then collapse to one observation per group \Adoref{4669}:
\begin{equation}
\boxed{\;\mathrm{in\_sum}_{k,d}(g) \;=\; \sum_{t}\sigma^{(d)}_{k,g,t}.\;}
\label{eq:insum}
\end{equation}
Stored as \code{in\_sum\_`k'\_`d'\_XX}.

\medskip\noindent
\textbf{What this is.} $\mathrm{in\_sum}_{k,d}(g)$ is group $g$'s contribution
to the score of the auxiliary regression of $\Delta Y$ on $\Delta X$ (with FEs)
for baseline cell $d$, after double-residualization. It is the
``control-side'' summary fed into $\mathrm{Den}_d^{-1}$ in Stage 3.

\subsection*{Step 2.5 \,---\, The actual residualization of $\Delta Y^{(\ell)}$ \Adoref{4679}}

\begin{equation}
\boxed{\;\Delta Y^{(\ell)}_{g,t}\ \leftarrow\ \Delta Y^{(\ell)}_{g,t} - \sum_{k=1}^{K}\thetad[k]\cdot \Delta X^{(\ell)}_{k,g,t}\quad\text{if }D_{g,1}=d.\;}
\label{eq:residY}
\end{equation}
In the code this is written as a loop over $k$ that subtracts one product at a
time. Effectively $\Delta Y^{(\ell)}_{g,t}$ is overwritten in place with the
controls-residualized long difference. After this line, the rest of the
event-study machinery (the part not specific to controls) treats
\code{diff\_y\_`l'\_XX} as if it were the raw outcome long-difference --
controls have vanished from the surface.

\medskip\noindent
\textbf{Important nuance \Adoref{4680}.} The replacement is unconditional on
$\Delta X^{(\ell)}_{k,g,t}$ being non-missing. If any control is missing,
$\Delta Y^{(\ell)}_{g,t}$ becomes missing and the cell is silently excluded
from the long-difference estimation. This is by design: it prevents the
estimator from using $(g,t)$ cells where the residualization is undefined.

% ============================================================================
\section[Stage 3 -- Variance correction U-var-X (L4905--4954)]{Stage 3 \,---\, The $U^{\mathrm{var},X}$ correction \Adoref{4905--4954}}

After Stage 2, the per-group variance score with no controls is, from L4904,
\begin{equation}
U^{(\pm,\mathrm{noX})}_{G,g,\ell}(t) \;=\; \1\!\bigl\{\ell\le T_g-1\bigr\}\cdot \frac{G}{N^{\pm}_{\ell}}\cdot \Bigl[\mathrm{dist}^{(\ell)}_{g,t} - \tfrac{N^{\pm,t}_{\ell,g}}{N^{\mathrm{c},(\ell)}_{g,t}}\mathrm{nc}^{(\ell)}_{g,t}\Bigr]\cdot \1\{\ell+1\le t\le T_g\}\cdot \Ngt\cdot \mathrm{DOF}^{(\ell)}_{g,t}\cdot \bigl(\Delta Y^{(\ell)}_{g,t} - \widehat{E}^{(\ell)}_{g,t}\bigr).
\label{eq:UnoX}
\end{equation}
With controls, Stage 3 adds a second term coming from the fact that
$\thetad$ is estimated. The block at L4908--4954 builds that second term,
\code{part2\_switch`+\!-'\_`l'\_XX}, group by group.

\subsection*{Step 3.1 \,---\, The bracketed score $\mathrm{br}_{d,j}$ \Adoref{4921--4938}}

For each baseline $d$ and each control coordinate $j$, the inner loop over $k$
\Adoref{4924--4935} computes
\begin{equation}
\bigl[\mathrm{Den}_d^{-1}\cdot \mathrm{in\_sum}_d\bigr]_{j}(g) \;=\; \sum_{k=1}^{K}\mathrm{Den}_d^{-1}[j,k]\cdot \mathrm{in\_sum}_{k,d}(g),
\label{eq:braprod}
\end{equation}
restricted to $\1\{D_{g,1}=d,\ F_g\ge 3\}$ \Adoref{4929}.
Then $\thetad[j]$ is subtracted \Adoref{4938}:
\begin{equation}
\boxed{\;\mathrm{br}_{d,j}(g) \;=\; \bigl[\mathrm{Den}_d^{-1}\cdot \mathrm{in\_sum}_d\bigr]_j(g) \,-\, \thetad[j].\;}
\label{eq:bra}
\end{equation}
Stored as \code{in\_brackets\_`d'\_`j'\_XX}.

\medskip\noindent
\textbf{What this is.} The asymptotic linearization of
$\thetad - \theta_d$ on the control sample is
$\thetad - \theta_d \approx \mathrm{Den}_d^{-1}\cdot \mathrm{in\_sum}_d / G + o_p(\cdot)$.
So $\mathrm{br}_{d,j}(g)$ is, up to sign, the influence-function coordinate of
$\thetad - \theta_d$ attributable to group $g$.

\subsection*{Step 3.2 \,---\, The per-group correction $\mathrm{part2}^{\pm}_{\ell}$ \Adoref{4941--4950}}

Multiply the influence function by the treatment-side moment and sum across
controls \Adoref{4942--4945}:
\begin{equation}
\mathrm{combined}^{\pm,d,\ell}(g) \;=\; \sum_{j=1}^{K} M^{\pm}_{d,\ell,j}\cdot \mathrm{br}_{d,j}(g),
\label{eq:combined}
\end{equation}
then sum over baseline cells \Adoref{4950}:
\begin{equation}
\boxed{\;\mathrm{part2}^{\pm}_{\ell}(g) \;=\; \sum_{d}\mathrm{combined}^{\pm,d,\ell}(g).\;}
\label{eq:part2}
\end{equation}
Stored as \code{part2\_switch`+\!-'\_`l'\_XX}.

\subsection*{Step 3.3 \,---\, Per-group variance score $U^{\pm,\mathrm{var},X}_{G,g,\ell}$ \Adoref{4957\,+\,}}

Sum the no-controls piece over $t$ within $g$ \Adoref{4957} and add the
controls correction:
\begin{equation}
U^{\pm,\mathrm{var},X}_{G,g,\ell} \;=\; \underbrace{\sum_{t} U^{(\pm,\mathrm{noX})}_{G,g,\ell}(t)}_{\text{\code{U\_Gg`l'\_var\_XX}}} \;+\; \mathrm{part2}^{\pm}_{\ell}(g).
\label{eq:UvarX}
\end{equation}

\subsection*{Step 3.4 \,---\, Asymptotic variance of the effect estimator}

Once $U^{\pm,\mathrm{var},X}_{G,g,\ell}$ is built for every group $g$, the
event-study variance of effect $\ell$ is the standard cross-group second
moment:
\begin{equation}
\widehat{\mathrm{Var}}\!\bigl(\widehat{\mathrm{Effect}}^{\pm}_{\ell}\bigr) \;=\; \frac{1}{G^{2}}\sum_{g=1}^{G} \bigl(U^{\pm,\mathrm{var},X}_{G,g,\ell}\bigr)^{2}
\;\;\stackrel{\text{cluster}(c)}{\longrightarrow}\;\;
\frac{1}{G^{2}}\sum_{c} \Bigl(\sum_{g\in c} U^{\pm,\mathrm{var},X}_{G,g,\ell}\Bigr)^{2},
\label{eq:Var}
\end{equation}
matching the cluster-robust convention implemented elsewhere in the file. The
$\mathrm{part2}^{\pm}_{\ell}(g)$ term is what makes \eqref{eq:Var} the
\emph{right} variance when controls are in the model; without it, the SEs
would understate the noise arising from $\thetad$ estimation.

% ============================================================================
\section[Visualizing the dependency graph]{Visualizing the dependency graph}

\begin{tcolorbox}[colback=codebg, colframe=navy, boxrule=0.6pt, arc=1pt]
\textbf{Where each Stage 1 output enters Stages 2--3:}\\[2pt]
\begin{tabular}{@{}l@{\;\;}c@{\;\;}l@{}}
$\thetad$  &$\;\to\;$& \textbf{Stage 2}, eq.~\eqref{eq:residY}, the residualization of $\Delta Y^{(\ell)}$.\\
           &$\;\to\;$& \textbf{Stage 3}, eq.~\eqref{eq:bra}, subtracted from the influence function.\\[3pt]
$\mathrm{Den}_d^{-1}$ &$\;\to\;$& \textbf{Stage 3}, eq.~\eqref{eq:braprod}, scaling of the score.\\[3pt]
$\mathrm{prod\_X}_{k,g,t}$ &$\;\to\;$& \textbf{Stage 2}, eq.~\eqref{eq:sigma}, the score summand.\\[3pt]
$\widehat{E}^{(d)}_{g,t}$ &$\;\to\;$& \textbf{Stage 2}, eq.~\eqref{eq:Ehat}/\eqref{eq:sigma}, the demean inside the score.\\
\end{tabular}
\end{tcolorbox}

% ============================================================================
\section[Cross-reference table (Stage 2)]{Cross-reference table \,---\, Stage 2}

\noindent\begin{tabular}{@{}lll@{}}
\toprule
Math symbol & \texttt{.ado} variable & Defined at \\
\midrule
$\Delta X^{(\ell)}_{k,g,t}$                & \code{diff\_X{k}\_`l'\_XX}           & L4593 \\
$\Delta X^{(\ell),N}_{k,g,t}$              & \code{diff\_X{k}\_`l'\_N\_XX}        & L4598 \\
$m^{\pm}_{g,d,\ell,k,t}$ (per-$t$)         & \code{m`+\!-'\_g\_`k'\_`d'\_`l'\_XX} & L4612 \\
$m^{\pm}_{g,d,\ell,k}$ (sum over $t$)      & \code{m`+\!-'\_`d'\_`k'\_`l'\_XX}    & L4615 \\
$M^{\pm}_{d,\ell,k}$ (avg over $g$)        & \code{M`+\!-'\_`d'\_`k'\_`l'\_XX}    & L4621 \\
$E^{\mathrm{denom}}_{d,\ell}(t)$           & \code{E\_hat\_denom\_`k'\_`d'\_XX}   & L4635/L4640 \\
$\widehat{E}^{(d,\ell)}_{g,t}$             & \code{E\_y\_hat\_gt\_`d'\_XX}        & L4647 \\
$T_d = \max_{h:D_{h,1}=d} F_h - 1$         & \code{T\_d\_XX}                      & L4579 \\
$N^{\mathrm{c}}_d$                         & \code{N\_c\_`d'\_XX}                 & L4659 \\
$\delta^{(d)}_{g,t}$ (DoF adjustment)      & \code{in\_sum\_temp\_adj\_`k'\_`d'\_XX} & L4662--4663 \\
$\sigma^{(d)}_{k,g,t}$ (score summand)     & \code{in\_sum\_temp\_`k'\_`d'\_XX}   & L4664 \\
$\mathrm{in\_sum}_{k,d}(g)$                & \code{in\_sum\_`k'\_`d'\_XX}         & L4669 \\
residualization $\Delta Y^{(\ell)} \mathrel{-}{=} \theta_d^{(k)}\Delta X^{(\ell)}_k$ & (overwrites \code{diff\_y\_`l'\_XX}) & L4679 \\
\bottomrule
\end{tabular}

\section[Cross-reference table (Stage 3)]{Cross-reference table \,---\, Stage 3}

\noindent\begin{tabular}{@{}lll@{}}
\toprule
Math symbol & \texttt{.ado} variable & Defined at \\
\midrule
$[\mathrm{Den}_d^{-1}\cdot \mathrm{in\_sum}_d]_{j,k}$ (term)  & \code{in\_brackets\_`d'\_`j'\_`k'\_temp\_XX} & L4929 \\
$[\mathrm{Den}_d^{-1}\cdot \mathrm{in\_sum}_d]_{j}$ (sum over $k$) & \code{in\_brackets\_`d'\_`j'\_XX} (incremented) & L4934 \\
$\mathrm{br}_{d,j}(g)$ (subtract $\thetad[j]$) & \code{in\_brackets\_`d'\_`j'\_XX} (final) & L4938 \\
$M^{\pm}_{d,\ell,j}\cdot \mathrm{br}_{d,j}(g)$ & \code{combined`+\!-'\_temp\_`d'\_`j'\_`l'\_XX} & L4942 \\
$\mathrm{combined}^{\pm,d,\ell}(g)$ (sum over $j$) & \code{combined`+\!-'\_temp\_`d'\_`l'\_XX} & L4945 \\
$\mathrm{part2}^{\pm}_{\ell}(g)$ (sum over $d$) & \code{part2\_switch`+\!-'\_`l'\_XX} & L4950 \\
$U^{(\pm,\mathrm{noX})}_{G,g,\ell}(t)$ summand & \code{U\_Gg`l'\_temp\_var\_XX} & L4904 \\
$U^{\pm,\mathrm{var},X}_{G,g,\ell}$ & \code{U\_Gg`l'\_var\_XX} (+ \code{part2\_switch`+\!-'\_`l'\_XX}) & L4957\,+\,$\dots$ \\
\bottomrule
\end{tabular}

% ============================================================================
\section[Placebos -- the symmetric story (L5145+)]{Placebos \,---\, the symmetric story \Adoref{5145+}}

The placebo path repeats Stages 2 and 3 with three substitutions:
\begin{itemize}[leftmargin=*, itemsep=2pt]
\item \textbf{Long differences are flipped in time.} \Adoref{5156}
      \[
      \Delta X^{(\ell),\mathrm{pl}}_{k,g,t} \;=\; X_{k,g,t-\ell} - X_{k,g,t-2\ell},
      \]
      so the placebo lives entirely \emph{before} the switch date.
\item \textbf{$\thetad$ is reused.} \Adoref{5177}
      The same \code{coefs\_sq\_`d'\_XX} computed in Stage 1 multiplies
      $\Delta X^{(\ell),\mathrm{pl}}_{k,g,t}$ to residualize
      \code{diff\_y\_pl\_`l'\_XX}. There is no separate ``placebo
      $\theta_d$''. This is by design: under parallel trends and
      no-anticipation, $\theta_d$ is identified pre-switch, exactly where
      placebos live.
\item \textbf{$\mathrm{Den}_d^{-1}$ and $\mathrm{in\_sum}_{k,d}$ are reused.}
      The placebo variance correction
      \code{part2\_pl\_switch`+\!-'\_`l'\_XX}
      multiplies the same influence function $\mathrm{br}_{d,j}$ by a placebo
      treatment-side moment $M^{\pm,\mathrm{pl}}_{d,\ell,j}$ \Adoref{5165--5172}.
\end{itemize}
The bookkeeping renames (\code{m\_pl}, \code{M\_pl}, \code{in\_sum\_pl},
\code{part2\_pl}) but the math is identical to \eqref{eq:bra}--\eqref{eq:part2}
with the placebo $\Delta X^{(\ell),\mathrm{pl}}$ replacing the effect
$\Delta X^{(\ell)}$.

% ============================================================================
\section[Why this design is elegant]{Why this design is elegant}

Three properties make this implementation robust:

\begin{enumerate}[leftmargin=*, itemsep=3pt]
\item \textbf{Stage 1 is computed once.} $\thetad$, $\mathrm{Den}_d^{-1}$,
      $\mathrm{prod\_X}_{k,g,t}$, $\widehat{E}^{(d)}_{g,t,\mathrm{int}}$
      depend only on $(g,t)$ data, not on the horizon $\ell$. They are
      recycled across every $\ell$ and across effects vs.\ placebos with no
      additional cost.
\item \textbf{Stage 2 is just an additive shift.} The residualization at
      L4679 overwrites $\Delta Y^{(\ell)}$ in place. After it, every
      downstream piece of code that has nothing to do with controls
      consumes the residualized long-difference as if it were the raw one.
      Adding HDFE (the previous note) requires no change to that
      downstream code.
\item \textbf{Stage 3 is a clean additive correction.} The
      $U^{\mathrm{var},X}$ term \eqref{eq:UvarX} is added to the no-controls
      variance score in one line at L4957$+$. Switching off controls
      (\code{`controls'==\!""}) skips L4908--4954 entirely and recovers the
      no-controls variance unchanged.
\end{enumerate}

The cost of this elegance is two nested loops over $K^{2}$ at L4921--4946,
which is fine for typical $K \le 5$ and dominates only when the user passes
many controls. If you ever add HDFE-derived controls (e.g.\ thousands of
dummies via \code{absorb()}), keep the absorbed dimensions out of $K$ -- that
is exactly what the HDFE plan does.

\end{document}
